% !TeX root = ../main.tex

% Edit this title for each week's entry.
\section{Week 7: Variance Decomposition Part 2}

\subsection{Review: Decomposition and $F$-Tests}

\begin{remark}[Sum of Squares Decomposition]
    \[\underbrace{\sum(y_i-\bar y)^2}_{SST}=\underbrace{\sum(\hat y_i-\bar y)^2}_{SS_{reg}}+\underbrace{\sum(y_i-\hat y_i)^2}_{RSS}\]
    \begin{itemize}
        \item $SST$: total variation in the response prior to fitting a model ($df=n-1$)
        \item $SS_{reg}$: variation explained by the model ($df=p$)
        \item $RSS$: variation left unexplained by the model ($df=n-p-1$)
    \end{itemize}
    A good model has $SS_{reg}>RSS$. But a model with more predictors (even non-significant
    ones) has smaller RSS:
    \[RSS_{small}\geq RSS_{big}\iff SS_{reg,small}\leq SS_{reg,big}\]
\end{remark}

\begin{remark}[ANOVA and Partial $F$-Tests]
    \begin{itemize}
        \item \textbf{ANOVA overall test}: $H_0$: all slopes are 0 vs.\ $H_A$: at least one slope is not 0,
        \[F^*=\frac{SS_{reg}/p}{RSS/(n-p-1)}\sim F(p,\,n-p-1)\text{ under }H_0\]
        \item \textbf{Partial $F$-test} (nested models): $H_0$: the full model is not better than
        the reduced model with $k$ predictors dropped,
        \[F^*=\frac{(RSS_{reduced}-RSS_{full})/k}{RSS_{full}/(n-p-1)}\sim F(k,\,n-p-1)\]
        \item The ANOVA overall test is the partial $F$-test with the intercept-only reduced model
        \aside{(\texttt{lm(y \textasciitilde{} 1)})}: there $RSS=SST$ and $SS_{reg}=0$; SST is the
        same in both models, so $RSS_{drop}=SST-RSS_{full}=SS_{reg,full}$.
    \end{itemize}
\end{remark}

\begin{proposition}[$t$-Tests and $F$-Tests Are Equivalent]
    \[F=\frac{\chi^2_{df_1}/df_1}{\chi^2_{df_2}/df_2}\sim F(df_1,df_2),\qquad
    T=\frac{Z}{\sqrt{\chi^2_{df}/df}}\sim t(df)\]
    If $Z\sim N(0,1)$ then $Z^2\sim\chi^2_1$, so $T^2\sim F(1,df)$: the $t$-test of
    $H_0:\beta_j=0$ is a partial $F$-test whose reduced model removes only $\beta_j$.
\end{proposition}

\heading{Module Outline}
\begin{itemize}
    \item Decomposition and measuring goodness
    \item Problems with related predictors
    \item Assessing and addressing multicollinearity
\end{itemize}

\subsection{Measuring Goodness}

\begin{remark}[Comparing Models]
    With the same dataset (same SST), we can compare models by which has the larger
    $SS_{reg}$. But suppose you and a friend work on separate problems on the same topic:
    different datasets, a response measuring the same characteristic (perhaps in different
    ways), and each a 2-predictor model with different predictors. Which model is better?
    \begin{itemize}
        \item We can't compare $SS_{reg}$, because each model's SST is different.
        \item We can't use the partial $F$-test, because the models are not nested.
        \item The ANOVA overall test only says whether a linear relationship exists; it doesn't
        quantify relative goodness.
    \end{itemize}
    $\Rightarrow$ We need a \textbf{unitless} measure.
\end{remark}

\begin{definition}[Coefficient of Determination $R^2$]
    Standardize the variation explained by SST:
    \[\boxed{R^2=\frac{SS_{reg}}{SST}=1-\frac{RSS}{SST}}\]
    \aside{Usually we use the $1-RSS/SST$ form.} $0\leq R^2\leq1$, and in SLR
    $R^2=\mathrm{corr}(X,Y)^2$. It is the \hlt{proportion of variation in the response
    explained by the model}.
\end{definition}

\begin{remark}[Notes about $R^2$]
    \begin{itemize}
        \item $R^2=1\iff RSS=0$: all points lie exactly on the fitted hyperplane, and
        $X_1,\dots,X_p$ account for all the variation in $Y$ (e.g.\ points exactly on a line with
        \texttt{lm(y\textasciitilde x)}, or a wiggly curve through every point with
        \texttt{lm(y\textasciitilde x+x\^{}2+...+x\^{}5)}).
        \item If $\hat\beta_1=\cdots=\hat\beta_p=0$, then $R^2=0$: the model explains no variation,
        i.e.\ there is no linear association between $X_1,\dots,X_p$ and $Y$ in the sample
        (a flat fit at $\bar y$).
        \item The closer $R^2$ is to 1, the greater the \textbf{degree of linear association}
        between $X_1,\dots,X_p$ and $Y$.
    \end{itemize}
\end{remark}

\begin{remark}[Interpreting $R^2$ with Care]
    \begin{itemize}
        \item There are no general rules about how large $R^2$ should be.
        \item $R^2$ can be large even when the linear regression model is not ideal.
        \item $R^2$ is not meaningful if a linear model is not appropriate.
        \item Very high $R^2$ is achievable by fitting complicated models with many predictors.
        \db{E.g.\ on the same curved data, \texttt{lm(y\textasciitilde x)} $\to$
        \texttt{lm(y\textasciitilde x+x\^{}2)} $\to$ \texttt{lm(y\textasciitilde x+...+x\^{}7)}: $R^2$
        keeps increasing, even as the last fit just chases noise.}
    \end{itemize}
\end{remark}

\begin{definition}[Adjusted Coefficient of Determination $R^2_{adj}$]
    Adding predictors \emph{always} decreases RSS, even if they're unrelated to the response,
    and we don't want more predictors than we need. Adjust with degrees of freedom (as in ANOVA):
    \[\boxed{R^2_{adj}=1-\frac{RSS/(n-p-1)}{SST/(n-1)}=1-\frac{s^2}{s_y^2}}\]
    \begin{itemize}
        \item This is \textbf{no longer} ``the proportion of variation explained by the model''.
        \item A bigger model is better only if $SS_{reg}$ has increased enough to compensate for
        the added complexity (the $-p$ in the denominator df penalizes extra predictors).
        \item $R^2_{adj}\leq R^2$, since $\frac{n-1}{n-p-1}\geq1$.
    \end{itemize}
\end{definition}

\begin{example}[Defects Data (Sheather p.\ 167)]
    Defect rates for 30 production lines: $Y=$ \texttt{Defective} (average \# defects per 1000
    parts), $X_1=$ \texttt{Temperature} of the process, $X_2=$ \texttt{Density} of the product,
    $X_3=$ \texttt{Rate} of production. The pairs plot shows $Y$ curving upward in Temperature and
    Rate and downward in Density, and the predictors strongly related to each other.
    \aside{Compare $y\sim x_1+x_2+x_3$ vs.\ $y\sim x_1+x_2+x_3+x_1^2+x_2^2+x_3^2$.}
\begin{verbatim}
defects <- read.table("defects.txt", header=T)
defects <- defects[,c(-1)]   # drop the first column, it's the case ID
n = nrow(defects)            # n = 30
model <- lm(Defective ~ ., data=defects)   # Y ~ . : Y on all other columns
SST <- var(defects$Defective)*(n-1)        # 10923.87
RSS <- sum(model$residuals^2)              # 1314.431
\end{verbatim}
    With $n=30$, $p=3$, $\Var(Y)=\frac{SST}{n-1}$ and $s^2=\frac{RSS}{n-p-1}$:
    \[s^2=50.555,\quad s_y^2=376.6853,\quad RSS=26(50.555)=1314.43,\quad SST=29(376.6853)=10923.87\]
    \[R^2=1-\frac{1314.43}{10923.87}\approx\dbb{0.88}\quad\text{(88\% of the variation in $Y$ explained by the model)}\]
    \[R^2_{adj}=1-\frac{s^2}{s_y^2}=1-\frac{50.555}{376.6853}\approx\dbb{0.866}\]
    Matching \texttt{summary(model)}:
\begin{verbatim}
            Estimate Std. Error t value Pr(>|t|)
(Intercept)  10.3244    65.9265   0.157   0.8768
Temperature  16.0779     8.2941   1.938   0.0635 .
Density      -1.8273     1.4971  -1.221   0.2332
Rate          0.1167     0.1306   0.894   0.3797
Residual standard error: 7.11 on 26 degrees of freedom
Multiple R-squared: 0.8797, Adjusted R-squared: 0.8658
F-statistic: 63.36 on 3 and 26 DF, p-value: 4.371e-12
\end{verbatim}
\end{example}

\begin{exercise}[Practice at Home]
    Find $R^2$ and $R^2_{adj}$ manually for various models on datasets we've used, check against
    the summary output, and interpret, e.g.\ \texttt{mpg \textasciitilde{} wt + hp} on \texttt{mtcars}.
\end{exercise}

\subsection{Problems with Related Predictors}

\begin{remark}[Recall: Conditional Mean Predictor Condition (Module 3)]
    To interpret residual plots we need each predictor's mean related to each other predictor
    in no more complicated way than linearly: $E(X_i\mid X_j)=\alpha_0+\alpha_1X_j$. Linear
    relationships or no relationship satisfy this (a random cloud, or a linear band, is fine);
    complicated relationships (e.g.\ a sine wave) can obfuscate residual plots.
\end{remark}

\begin{remark}[Related Predictors and Correlation]
    \begin{itemize}
        \item Ideally there is no relationship between predictors --- this almost never happens
        unless an experiment is carefully designed for it.
        \item A linear trend between predictors means they are \textbf{correlated} or
        \textbf{collinear}.
        \item Recall: the correlation $\rho$ measures the strength of linear association between
        two continuous variables.
        \item From the house pricing example (\texttt{price \textasciitilde{} \#bedrooms + \#bathrooms},
        \#bedrooms $\approx$ \#bathrooms): the more positively correlated $X_1$ and $X_2$, the more
        negatively correlated $\hat\beta_1$ and $\hat\beta_2$ (and vice versa), \textbf{and} the higher
        the variances of $\hat\beta_1$ and $\hat\beta_2$.
    \end{itemize}
\end{remark}

\begin{example}[House Pricing: Unstable Estimates]
    $Y=1+3X_1+2X_2+\epsilon$, $\epsilon\overset{iid}\sim N(0,3)$, $X_1\approx X_2$.
    \aside{Fits like $\hat Y=1+5X_1$, $\hat Y=1+4X_1+X_2$, $\hat Y=1+5X_2$, \dots\ are all
    plausible.}
    \begin{center}
        \begin{tabular}{lrrrr}
            \toprule
            & \multicolumn{2}{c}{\texttt{set.seed(1234)}} & \multicolumn{2}{c}{different seed} \\
            & Estimate (SE) & p-value & Estimate (SE) & p-value \\
            \midrule
            \texttt{numBedrooms} $\hat\beta_1$ & 2.3661 (2.2732) & 0.301 & 4.4332 (2.1403) & 0.041 \\
            \texttt{numBathrooms} $\hat\beta_2$ & 2.4111 (2.2772) & 0.292 & 0.7338 (2.1164) & 0.730 \\
            \bottomrule
        \end{tabular}
    \end{center}
    The data-generating process is the same, but the estimates are very different, the variances
    are large, and \textbf{the significance changed as well}. \texttt{vcov(model)} gives
    $\widehat{\Var}(\hat{\boldsymbol\beta})=\hat\sigma^2(\mathbf{X}'\mathbf{X})^{-1}$; for the first seed,
    \[\mathrm{Corr}(\hat\beta_1,\hat\beta_2)=\frac{\Cov(\hat\beta_1,\hat\beta_2)}{\sqrt{\Var(\hat\beta_1)\Var(\hat\beta_2)}}
    =\frac{-5.16}{\sqrt{5.17\cdot5.19}}\approx-1\]
    \db{(Second seed: $-4.513/\sqrt{4.581\cdot4.479}\approx-1$ too.)}
\end{example}

\begin{definition}[Rank of the Design Matrix]
    Perfect correlation ($\pm1$) between predictors --- one perfectly linearly related to another
    --- causes issues. More generally, if one column of $\mathbf{X}$ is a \textbf{linear
    combination of the other columns}, the predictors are dependent, $\mathbf{X}$ is not of full
    column rank, and $\mathbf{X}'\mathbf{X}$ is \textbf{not invertible}: we can't estimate the
    linear relationship.
\end{definition}

\begin{example}[Matrix Not Invertible]
    $n=5$, $p=3$:
    \[\mathbf{X}=\begin{pmatrix}1&1&4&3\\1&2&2&3\\1&3&8&7\\1&2&2&3\\1&3&8&7\end{pmatrix},\qquad
    \mathbf{X}'\mathbf{X}=\begin{pmatrix}5&11&24&23\\11&27&60&57\\24&60&152&136\\23&57&136&125\end{pmatrix}\]
    Here $X_3=X_1+0.5X_2$, so $\mathbf{X}'\mathbf{X}$ has rank 3 ($C_4=C_2+0.5C_3$): not
    invertible. \texttt{solve(t(X)\%*\%X)} errors: ``system is computationally singular''.
    \\ Yet no pair of predictors is perfectly correlated --- \texttt{cor(X)} (the correlation matrix
    of $\mathbf{X}$) gives $r(X_1,X_2)=0.709$, $r(X_1,X_3)=0.873$, $r(X_2,X_3)=0.963$. Pairwise
    correlations tell us whether any 2 predictors are related, but we also need to know whether a
    predictor is a linear combination of \emph{several} others.
\end{example}

\begin{definition}[Multicollinearity]
    When one predictor can be modelled really well using the others. Its impact on the model
    depends on how strong the relationship between the predictors is, and how much of the
    variation in $Y$ these predictors explain.
\end{definition}

\begin{remark}[Impact of Multicollinearity]
    \begin{itemize}
        \item Instability in how predictors are used to estimate the mean responses: the model
        can't distinguish how much variation is explained only by $X_1$ vs.\ only by $X_2$, so
        \hlt{$\Var(\hat{\boldsymbol\beta})$ is large}.
        \item House prices: hard to separate what's explained by \#bedrooms vs.\ \#bathrooms. For
        a house with many more bedrooms than bathrooms, the prediction can change a lot from small
        differences in the training data --- \textbf{extrapolation is very dangerous here}.
        \item It prevents us from properly understanding the predictor--response relationship:
        \begin{itemize}
            \item \textbf{Wrong estimated coefficients}: signs may be wrong compared to the literature.
            \item \textbf{Contradictory significance}: many predictors insignificant while the
            overall $F$-test is highly significant.
            \item \textbf{Inflated variances}: standard errors much larger than they should be.
        \end{itemize}
        All of these appear in the simulated housing example!
    \end{itemize}
\end{remark}

\begin{remark}[Structural Multicollinearity]
    Sometimes we create multicollinearity on purpose, as a necessary result of the model
    structure we choose to satisfy the assumptions:
    \begin{itemize}
        \item \textbf{Polynomial regression} for linearity, e.g.\ $\hat y_i=\hat\beta_0+\hat\beta_1x_i+\hat\beta_2x_i^2$,
        so column $C_3$ of $\mathbf{X}$ is $(C_2)^2$.
        \item \textbf{Interaction terms with their main effects}: if you include an interaction you
        generally also include the main effects, e.g.\
        $\hat y_i=\hat\beta_0+\hat\beta_1x_{i1}+\hat\beta_2x_{i2}+\hat\beta_3x_{i1}x_{i2}$.
    \end{itemize}
\end{remark}

\subsection{Assessing and Addressing Multicollinearity}

\begin{remark}[EDA and Preliminary Assessment]
    A preliminary assessment can help justify decisions about handling multicollinearity.
    \begin{itemize}
        \item \textbf{Pairwise correlations} (correlation matrix of $X$): two predictors at a time.
        Can be done early, but may not show the full extent of the problem (recall the
        non-invertible example) and may miss non-linear relationships.
        \item \textbf{Pairs plots}: two predictors at a time; can reveal non-linear relationships
        (not necessarily multicollinearity --- e.g.\ an additional condition 2 issue from Module 3).
    \end{itemize}
    But we need to see how \emph{multiple} predictors act together to explain the response.
\end{remark}

\begin{definition}[Variance Inflation Factor (VIF), 2 Predictors]
    The VIF quantifies how much larger the variance of a coefficient is due to
    multicollinearity. For $y_i=\beta_0+\beta_1x_{i1}+\beta_2x_{i2}+\varepsilon_i$, with
    $\Var(\hat{\boldsymbol\beta})=\sigma^2(\mathbf{X}'\mathbf{X})^{-1}$, it can be shown that
    \[\boxed{\Var(\hat\beta_j)=\underbrace{\frac{1}{1-r_{12}^2}}_{VIF}\times\frac{\sigma^2}{(n-1)s_{x_j}^2}},\qquad j=1,2\]
    where $r_{12}$ is the sample correlation between $x_1$ and $x_2$ and $s_{x_j}$ the sample SD of $x_j$.
    \begin{itemize}
        \item $\frac{1}{1-r_{12}^2}$ is how inflated $\Var(\hat\beta_j)$ is due to correlation
        between the predictors.
        \item If the predictors are independent, $\Var(\hat\beta_j)$ is close to $\frac{\sigma^2}{(n-1)s_{x_j}^2}$.
        \item As $|r_{12}|$ increases, the VIF increases.
    \end{itemize}
\end{definition}

\begin{example}[VIF for House Pricing]
    Computing the VIF by hand and with \texttt{car::vif()}:
\begin{verbatim}
r <- cor(house1$numBedrooms, house1$numBathrooms)
1/(1-r^2)        # 160.1022  (second simulation: 135.8856)
car::vif(model2) # numBedrooms 160.1022, numBathrooms 160.1022
\end{verbatim}
    Both simulations have huge VIFs. \aside{With only 2 predictors the VIFs for $\hat\beta_1$ and
    $\hat\beta_2$ are the same.}
\end{example}

\begin{definition}[General VIF]
    For more than 2 predictors, define $R_j^2$: the proportion of the variation in $X_j$ explained
    by the \textbf{remaining predictors}, i.e.\ the $R^2$ from regressing $X_j$ on the rest.
    E.g.\ for $y\sim x_1+x_2+x_3$: $x_1\sim x_2+x_3\to R_1^2$, $x_2\sim x_1+x_3\to R_2^2$,
    $x_3\sim x_1+x_2\to R_3^2$. Then
    \[\boxed{\Var(\hat\beta_j)=\frac{1}{1-R_j^2}\times\frac{\sigma^2}{(n-1)s_{x_j}^2},\qquad VIF_j=\frac{1}{1-R_j^2}},\quad j=1,\dots,p\]
    We replaced $r_{12}^2$ with $R_j^2$, which tells us how much of the information in $X_j$ is
    already provided by the other predictors.
\end{definition}

\begin{remark}[Notes about the VIF]
    \begin{itemize}
        \item Generalizes to models of any size.
        \item Can only be calculated for \textbf{numerical} predictors ($X_j$ is the response in
        $X_j\sim$ rest of the $X$'s).
        \item The more of the variation in $X_j$ the remaining $p-1$ predictors explain, the larger
        the inflation.
        \item Rule of thumb for this class: \[\boxed{VIF>5\ \Rightarrow\ \text{severe multicollinearity}}\]
        a sign the variances of our estimates are very inflated; watch out for potentially
        incorrect interpretations.
        \item Ensure the model assumptions hold before interpreting the VIF.
    \end{itemize}
\end{remark}

\begin{example}[VIF for the Defects Data]
    For \texttt{lm(Defective \textasciitilde{} Temperature + Density + Rate)}, the variance of the
    \texttt{Temperature} coefficient (a numeric predictor) depends on how much Temperature overlaps
    with Density and Rate: fit $X_1\sim X_2+X_3$.
\begin{verbatim}
model1a <- lm(Temperature ~ Density + Rate, data=defects)
summary(model1a)   # Multiple R-squared: 0.9255
\end{verbatim}
    $R_1^2=0.9255$: 92.55\% of the variation in $X_1$ is explained by $X_2$ and $X_3$, so
    \[VIF_1=\frac1{1-R_1^2}=\frac1{1-0.9255}=\dbb{13.42}\]
\begin{verbatim}
library(car)
vif(model1)   # fitted model object  lm(y ~ x1 + x2 + x3)
# Temperature   Density      Rate
#   13.431614 14.508872  6.642619
\end{verbatim}
    $\Var(\hat\beta_1)$, $\Var(\hat\beta_2)$, $\Var(\hat\beta_3)$ are inflated by factors of 13.4,
    14.5 and 6.6 --- all above 5, so there is significant multicollinearity among all predictors.
    Unsurprising: the pairs plot showed strong relationships between the predictors.
    \aside{With more than 2 predictors the VIFs differ across predictors.}
\end{example}

\begin{remark}[Addressing Multicollinearity]
    Two main ways:
    \begin{itemize}
        \item \textbf{Collect more data} (not usually feasible).
        \item \textbf{Respecify the model} (what's typically done): remove some predictors or
        change their form.
    \end{itemize}
    Multicollinearity means different predictors explain the same parts of the overall variation
    in $Y$ (the SST); removing one or some predictors removes some of the overlap. Think carefully
    about which to remove:
    \begin{itemize}
        \item keep predictors related to the research question;
        \item consider how model goodness is impacted (assumptions, statistical significance,
        width of prediction intervals);
        \item removing predictors must be justified (variable selection, Module 9).
    \end{itemize}
\end{remark}

\begin{remark}[Module Takeaways]
    \begin{itemize}
        \item How are $R^2$ and $R^2_{adj}$ computed and interpreted, and why do we need both?
        \db{$R^2$: proportion of variation in $Y$ explained by $X$; $R^2_{adj}$: adjusts for the
        size of the model.}
        \item What does the VIF measure, and why is it useful for detecting multicollinearity?
        \item How do we identify multicollinearity, and why is it important to identify?
    \end{itemize}
\end{remark}
